\chapter{Lógica matemática y teoría elemental de conjuntos}

\section{Curso de un mes}

\begin{lesson}[id=curso-guia-general,weight=0]

Duración: 30 días. Cada día contiene diez actividades: seis de respuesta corta, dos de desarrollo y dos de opción múltiple. Las entregas indicadas se realizan en los días señalados. Trabaja en orden y conserva la notación utilizada en cada ejercicio.

\end{lesson}

\begin{lesson}[id=convenciones-notacion,weight=0]

Convenciones: \(V\) significa verdadero y \(F\) falso; \(\neg\) negación; \(\land\) conjunción; \(\lor\) disyunción; \(\to\) condicional; \(\in\) pertenencia; \(\subseteq\) inclusión; \(\varnothing\) conjunto vacío; \(P(A)\) conjunto potencia. En los ejercicios de conjuntos, el universo se indicará cuando sea necesario. Para los naturales, se adopta aquí \(\mathbb{N}=\{0,1,2,\ldots\}\).

\end{lesson}

\section{Día 1: Enunciados y proposiciones}

\begin{lesson}[id=dia-01-lectura,weight=0]

Distinguir enunciados declarativos de preguntas, órdenes y expresiones abiertas.

\end{lesson}

\begin{activity}[id=dia-01-corta-01,type=choice,weight=1]

Clasifica como proposición o no proposición: «7 es primo». Justifica en una frase.
\option{a}{Es proposición}
\option{b}{No es proposición}
\correct{b}


\end{activity}


\begin{activity}[id=dia-01-corta-02,type=short,weight=1]

Clasifica: «¿Cuánto vale x?»

\end{activity}

\begin{activity}[id=dia-01-corta-03,type=short,weight=1]

Clasifica: «Cierra la puerta».

\end{activity}

\begin{activity}[id=dia-01-corta-04,type=short,weight=1]

Indica si «La Tierra gira alrededor del Sol» tiene valor de verdad.

\end{activity}

\begin{activity}[id=dia-01-corta-05,type=short,weight=1]

Escribe dos proposiciones verdaderas y dos falsas.

\end{activity}

\begin{activity}[id=dia-01-corta-06,type=short,weight=1]

Convierte la frase «Hay un número natural mayor que 10» en una proposición simbólica o precisa.

\end{activity}

\begin{activity}[id=dia-01-desarrollo-01,type=long,weight=2]

Explica por qué una pregunta no es una proposición en lógica clásica. Incluye dos ejemplos y contrasta con dos proposiciones.

\end{activity}

\begin{activity}[id=dia-01-desarrollo-02,type=long,weight=2]

Escribe cinco enunciados cotidianos y analiza cuáles tienen un valor de verdad definido. Corrige los que no sean proposiciones para convertirlos en proposiciones.

\end{activity}

\begin{activity}[id=dia-01-opcion-01,type=choice,weight=1]

¿Cuál es una proposición?

\option{a}{Una pregunta}

\option{b}{5 es impar}

\option{c}{¡Ven aquí!}

\correct{b}

\end{activity}

\begin{activity}[id=dia-01-opcion-02,type=choice,weight=1]

¿Cuál de estos enunciados tiene valor de verdad?

\option{a}{¿Qué hora es?}

\option{b}{\(x+2=5\), sin fijar \(x\)}

\option{c}{11 es mayor que 4}

\correct{c}

\end{activity}

\section{Día 2: Proposiciones simples y compuestas}

\begin{lesson}[id=dia-02-lectura,weight=0]

Identificar proposiciones simples y reconocer conectivos lógicos.

\end{lesson}

\begin{activity}[id=dia-02-corta-01,type=short,weight=1]

Separa las proposiciones simples de: «Llueve y hace frío».

\end{activity}

\begin{activity}[id=dia-02-corta-02,type=short,weight=1]

Simboliza: «Estudio y trabajo», usando \(p\) y \(q\).

\end{activity}

\begin{activity}[id=dia-02-corta-03,type=short,weight=1]

Simboliza: «No llueve», usando \(p\) para «llueve».

\end{activity}

\begin{activity}[id=dia-02-corta-04,type=short,weight=1]

Escribe una proposición compuesta con «o».

\end{activity}

\begin{activity}[id=dia-02-corta-05,type=short,weight=1]

Identifica el conectivo principal de: \(\neg\)p \(\land\) q.

\end{activity}

\begin{activity}[id=dia-02-corta-06,type=short,weight=1]

Traduce a lenguaje cotidiano: \(p \to q\).

\end{activity}

\begin{activity}[id=dia-02-desarrollo-01,type=long,weight=2]

Construye tres proposiciones compuestas distintas a partir de dos proposiciones simples. Señala en cada una el conectivo principal.

\end{activity}

\begin{activity}[id=dia-02-desarrollo-02,type=long,weight=2]

Explica la diferencia entre una proposición simple y una compuesta y analiza la estructura de dos ejemplos propios.

\end{activity}

\begin{activity}[id=dia-02-opcion-01,type=choice,weight=1]

El símbolo \(\land\) se lee:

\option{a}{O}

\option{b}{Y}

\option{c}{Si... entonces}

\correct{b}

\end{activity}

\begin{activity}[id=dia-02-opcion-02,type=choice,weight=1]

La expresión \(\neg p\) representa:

\option{a}{La negación de \(p\)}

\option{b}{La conjunción de \(p\) y \(q\)}

\option{c}{La implicación de \(p\) en \(q\)}

\correct{a}

\end{activity}

\section{Día 3: Negación}

\begin{lesson}[id=dia-03-lectura,weight=0]

Construir negaciones correctas y distinguir negación gramatical de negación lógica.

\end{lesson}

\begin{activity}[id=dia-03-corta-01,type=short,weight=1]

Escribe \(\neg p\) si \(p\): «8 es par».

\end{activity}

\begin{activity}[id=dia-03-corta-02,type=short,weight=1]

Niega: «El número 9 es primo».

\end{activity}

\begin{activity}[id=dia-03-corta-03,type=short,weight=1]

Niega: «Todos los perros son negros».

\end{activity}

\begin{activity}[id=dia-03-corta-04,type=short,weight=1]

Niega: «Existe un número natural mayor que 100».

\end{activity}

\begin{activity}[id=dia-03-corta-05,type=short,weight=1]

Escribe la doble negación de \(p\).

\end{activity}

\begin{activity}[id=dia-03-corta-06,type=short,weight=1]

Determina el valor de verdad de \(\neg p\) cuando \(p\) es falsa.

\end{activity}

\begin{activity}[id=dia-03-desarrollo-01,type=long,weight=2]

Explica por qué la negación de «Todos los estudiantes aprobaron» no es «Ningún estudiante aprobó». Escribe la negación correcta.

\end{activity}

\begin{activity}[id=dia-03-desarrollo-02,type=long,weight=2]

Formula cuatro proposiciones cuantificadas informalmente y escribe sus negaciones usando «existe» o «para todo».

\end{activity}

\begin{activity}[id=dia-03-opcion-01,type=choice,weight=1]

Si p es verdadera, \(\neg p\) es:

\option{a}{Verdadera}

\option{b}{Falsa}

\option{c}{Indeterminada}

\correct{b}

\end{activity}

\begin{activity}[id=dia-03-opcion-02,type=choice,weight=1]

La negación de «Todos tienen libro» es:

\option{a}{Nadie tiene libro}

\option{b}{Alguien no tiene libro}

\option{c}{Todos no tienen libro}

\correct{b}

\end{activity}

\section{Día 4: Conjunción, disyunción y condicional}

\begin{lesson}[id=dia-04-lectura,weight=0]

Interpretar \(\land\), \(\lor\) y \(\to\) en lenguaje natural y simbólico.

\end{lesson}

\begin{activity}[id=dia-04-corta-01,type=short,weight=1]

Simboliza «\(p\) y \(q\)».

\end{activity}

\begin{activity}[id=dia-04-corta-02,type=short,weight=1]

Simboliza «\(p\) o \(q\)».

\end{activity}

\begin{activity}[id=dia-04-corta-03,type=short,weight=1]

Simboliza «si \(p\), entonces \(q\)».

\end{activity}

\begin{activity}[id=dia-04-corta-04,type=short,weight=1]

Traduce \(p\) \(\lor\) \(q\) a palabras.

\end{activity}

\begin{activity}[id=dia-04-corta-05,type=short,weight=1]

Escribe un ejemplo cotidiano de \(p\) \(\to\) \(q\).

\end{activity}

\begin{activity}[id=dia-04-corta-06,type=short,weight=1]

Identifica el conectivo principal de (\(p\) \(\land\) \(q\)) \(\to\) \(r\).

\end{activity}

\begin{activity}[id=dia-04-desarrollo-01,type=long,weight=2]

Explica la diferencia entre «o» inclusiva en lógica y el uso exclusivo que a veces tiene en el lenguaje cotidiano.

\end{activity}

\begin{activity}[id=dia-04-desarrollo-02,type=long,weight=2]

Crea una situación con dos proposiciones simples y escribe su conjunción, disyunción y condicional. Explica cada expresión.

\end{activity}

\begin{activity}[id=dia-04-opcion-01,type=choice,weight=1]

El condicional \(p\) \(\to\) \(q\) es falso únicamente cuando:

\option{a}{\(p\) es falsa y \(q\) verdadera}

\option{b}{\(p\) es verdadera y \(q\) falsa}

\option{c}{ambas son verdaderas}

\correct{b}

\end{activity}

\begin{activity}[id=dia-04-opcion-02,type=choice,weight=1]

La disyunción \(p\) \(\lor\) \(q\) es verdadera cuando:

\option{a}{Al menos una es verdadera}

\option{b}{Ambas son falsas}

\option{c}{Solo \(p\) es falsa}

\correct{a}

\end{activity}

\section{Día 5: Tablas de verdad: negación y conectivos}

\begin{lesson}[id=dia-05-lectura,weight=0]

Construir tablas de verdad básicas.

\end{lesson}

\begin{activity}[id=dia-05-corta-01,type=short,weight=1]

Completa la columna \(\neg p\) para \(p=V,F\).

\end{activity}

\begin{activity}[id=dia-05-corta-02,type=short,weight=1]

Completa \(p\) \(\land\) \(q\) para (\(V\),\(V\)).

\end{activity}

\begin{activity}[id=dia-05-corta-03,type=short,weight=1]

Completa \(p\) \(\land\) \(q\) para (\(V\),\(F\)).

\end{activity}

\begin{activity}[id=dia-05-corta-04,type=short,weight=1]

Completa \(p\) \(\lor\) \(q\) para (\(F\),\(F\)).

\end{activity}

\begin{activity}[id=dia-05-corta-05,type=short,weight=1]

Completa \(p\) \(\to\) \(q\) para (\(F\),\(V\)).

\end{activity}

\begin{activity}[id=dia-05-corta-06,type=short,weight=1]

¿Cuántas filas tiene una tabla con dos variables proposicionales?

\end{activity}

\begin{activity}[id=dia-05-desarrollo-01,type=long,weight=2]

Construye las tablas completas de \(p\) \(\land\) \(q\) y \(p\) \(\lor\) \(q\). Explica qué filas determinan sus valores.

\end{activity}

\begin{activity}[id=dia-05-desarrollo-02,type=long,weight=2]

Construye la tabla de \(p\) \(\to\) \(q\) y explica por qué solo una combinación la hace falsa.

\end{activity}

\begin{activity}[id=dia-05-opcion-01,type=choice,weight=1]

Con dos variables hay:

\option{a}{2 filas}

\option{b}{3 filas}

\option{c}{4 filas}

\correct{b}

\end{activity}

\begin{activity}[id=dia-05-opcion-02,type=choice,weight=1]

\(\neg\)F equivale a:

\option{a}{F}

\option{b}{V}

\option{c}{p}

\correct{b}

\end{activity}

\begin{deliverable}[id=entrega-dia-05,type=upload,weight=3]

Entrega un PDF con las tablas de verdad de la negación, conjunción, disyunción y condicional. Incluye una explicación breve de cada columna final.

\end{deliverable}

\section{Día 6: Tablas de verdad compuestas}

\begin{lesson}[id=dia-06-lectura,weight=0]

Evaluar expresiones con varios conectivos.

\end{lesson}

\begin{activity}[id=dia-06-corta-01,type=short,weight=1]

Construye la tabla de \(\neg p\) \(\land\) \(q\).

\end{activity}

\begin{activity}[id=dia-06-corta-02,type=short,weight=1]

Evalúa \(p\) \(\lor\) \(\neg q\) cuando \(p=F\) y \(q=F\).

\end{activity}

\begin{activity}[id=dia-06-corta-03,type=short,weight=1]

Evalúa (\(p\) \(\land\) \(q\)) \(\to\) \(p\) cuando \(p=V\) y \(q=F\).

\end{activity}

\begin{activity}[id=dia-06-corta-04,type=short,weight=1]

Identifica el conectivo principal de \(\neg p\) \(\lor\) \(q\).

\end{activity}

\begin{activity}[id=dia-06-corta-05,type=short,weight=1]

¿Cuántas filas requiere una tabla de tres variables?

\end{activity}

\begin{activity}[id=dia-06-corta-06,type=short,weight=1]

Evalúa \(p\) \(\to\) (q \(\land\) p) para \(p=F\) y \(q=V\).

\end{activity}

\begin{activity}[id=dia-06-desarrollo-01,type=long,weight=2]

Construye la tabla completa de (\(p\) \(\land\) \(q\)) \(\to\) \(r\). Muestra el orden de evaluación de los conectivos.

\end{activity}

\begin{activity}[id=dia-06-desarrollo-02,type=long,weight=2]

Compara \(\neg\)(\(p\) \(\land\) \(q\)) con (\(\neg p\) \(\land\) \(\neg q\)) mediante tablas. Indica si tienen siempre el mismo valor.

\end{activity}

\begin{activity}[id=dia-06-opcion-01,type=choice,weight=1]

Una tabla con tres variables tiene:

\option{a}{4 filas}

\option{b}{6 filas}

\option{c}{8 filas}

\correct{c}

\end{activity}

\begin{activity}[id=dia-06-opcion-02,type=choice,weight=1]

El conectivo principal de (\(p\) \(\lor\) \(q\)) \(\land\) \(r\) es:

\option{a}{\(\lor\)}

\option{b}{\(\land\)}

\option{c}{\(\neg\)}

\correct{b}

\end{activity}

\section{Día 7: Tautología, contradicción y contingencia}

\begin{lesson}[id=dia-07-lectura,weight=0]

Clasificar fórmulas mediante sus valores de verdad.

\end{lesson}

\begin{activity}[id=dia-07-corta-01,type=short,weight=1]

Define tautología en una frase.

\end{activity}

\begin{activity}[id=dia-07-corta-02,type=short,weight=1]

Define contradicción en una frase.

\end{activity}

\begin{activity}[id=dia-07-corta-03,type=short,weight=1]

Define contingencia en una frase.

\end{activity}

\begin{activity}[id=dia-07-corta-04,type=short,weight=1]

Clasifica \(p\) \(\lor\) \(\neg p\).

\end{activity}

\begin{activity}[id=dia-07-corta-05,type=short,weight=1]

Clasifica \(p\) \(\land\) \(\neg p\).

\end{activity}

\begin{activity}[id=dia-07-corta-06,type=short,weight=1]

¿Qué indica una columna final con valores \(V\) y \(F\)?

\end{activity}

\begin{activity}[id=dia-07-desarrollo-01,type=long,weight=2]

Construye una tabla para \(p\) \(\to\) \(p\) y clasifícala. Justifica la clasificación.

\end{activity}

\begin{activity}[id=dia-07-desarrollo-02,type=long,weight=2]

Da un ejemplo de tautología, contradicción y contingencia, y verifica cada uno con una tabla.

\end{activity}

\begin{activity}[id=dia-07-opcion-01,type=choice,weight=1]

p\(\lor\)\(\neg\)p es una:

\option{a}{Tautología}

\option{b}{Contradicción}

\option{c}{Contingencia}

\correct{a}

\end{activity}

\begin{activity}[id=dia-07-opcion-02,type=choice,weight=1]

Una contradicción es falsa:

\option{a}{En algunas filas}

\option{b}{En todas las filas}

\option{c}{En ninguna fila}

\correct{b}

\end{activity}

\section{Día 8: Equivalencias lógicas}

\begin{lesson}[id=dia-08-lectura,weight=0]

Reconocer equivalencias y aplicar leyes básicas.

\end{lesson}

\begin{activity}[id=dia-08-corta-01,type=short,weight=1]

Escribe la ley de doble negación.

\end{activity}

\begin{activity}[id=dia-08-corta-02,type=short,weight=1]

Completa: \(\neg\)(\(\neg p\)) \(\equiv\) __.

\end{activity}

\begin{activity}[id=dia-08-corta-03,type=short,weight=1]

¿Qué expresa p \(\land\) q \(\equiv\) q \(\land\) p?

\end{activity}

\begin{activity}[id=dia-08-corta-04,type=short,weight=1]

Aplica De Morgan a \(\neg\)(\(p\) \(\land\) \(q\)).

\end{activity}

\begin{activity}[id=dia-08-corta-05,type=short,weight=1]

Aplica De Morgan a \(\neg\)(\(p\) \(\lor\) \(q\)).

\end{activity}

\begin{activity}[id=dia-08-corta-06,type=short,weight=1]

Indica si \(p\) \(\to\) \(q\) equivale a \(\neg\)p \(\lor\) \(q\).

\end{activity}

\begin{activity}[id=dia-08-desarrollo-01,type=long,weight=2]

Verifica mediante tabla una ley de De Morgan.

\end{activity}

\begin{activity}[id=dia-08-desarrollo-02,type=long,weight=2]

Explica qué significa que dos fórmulas sean lógicamente equivalentes y presenta un ejemplo verificado.

\end{activity}

\begin{activity}[id=dia-08-opcion-01,type=choice,weight=1]

La doble negación establece:

\option{a}{\(\neg \neg p\) \(\equiv  p\)}

\option{b}{\(\neg p \equiv p\)}

\option{c}{p\(\land\)p \(\equiv\) \(\neg\)p}

\correct{a}

\end{activity}

\begin{activity}[id=dia-08-opcion-02,type=choice,weight=1]

\(\neg\)(\(p\) \(\land\) \(q\)) equivale a:

\option{a}{\(\neg p\) \(\land\) \(\neg q\)}

\option{b}{\(\neg p\) \(\lor\) \(\neg q\)}
\option{c}{\(p\) \(\lor\) \(q\)}

\correct{b}

\end{activity}

\section{Día 9: Argumentos e inferencias}

\begin{lesson}[id=dia-09-lectura,weight=0]

Distinguir premisas, conclusión y validez de un argumento.

\end{lesson}

\begin{activity}[id=dia-09-corta-01,type=short,weight=1]

Identifica las premisas de: Si llueve, la calle se moja; llueve; por tanto, la calle se moja.

\end{activity}

\begin{activity}[id=dia-09-corta-02,type=short,weight=1]

Identifica la conclusión del argumento anterior.

\end{activity}

\begin{activity}[id=dia-09-corta-03,type=short,weight=1]

Escribe un argumento con dos premisas y una conclusión.

\end{activity}

\begin{activity}[id=dia-09-corta-04,type=short,weight=1]

¿La validez depende de la verdad efectiva de las premisas o de la forma?

\end{activity}

\begin{activity}[id=dia-09-corta-05,type=short,weight=1]

Separa premisa y conclusión usando «por tanto».

\end{activity}

\begin{activity}[id=dia-09-corta-06,type=short,weight=1]

Representa con p y q: «Si p, entonces q; p; luego q».

\end{activity}

\begin{activity}[id=dia-09-desarrollo-01,type=long,weight=2]

Analiza un argumento cotidiano: enumera sus premisas, conclusión y forma lógica.

\end{activity}

\begin{activity}[id=dia-09-desarrollo-02,type=long,weight=2]

Explica la diferencia entre un argumento válido y uno con premisas verdaderas. Incluye un ejemplo de cada idea.

\end{activity}

\begin{activity}[id=dia-09-opcion-01,type=choice,weight=1]

En un argumento, la conclusión es:

\option{a}{La afirmación que se pretende derivar}

\option{b}{Cualquier ejemplo}

\option{c}{Una pregunta}

\correct{a}

\end{activity}

\begin{activity}[id=dia-09-opcion-02,type=choice,weight=1]

La validez lógica depende principalmente de:

\option{a}{La popularidad del argumento}

\option{b}{Su estructura inferencial}

\option{c}{La longitud}

\correct{b}

\end{activity}

\section{Día 10: Reglas de inferencia I}

\begin{lesson}[id=dia-10-lectura,weight=0]

Aplicar modus ponens y modus tollens.

\end{lesson}

\begin{activity}[id=dia-10-corta-01,type=short,weight=1]

Escribe el esquema de modus ponens.

\end{activity}

\begin{activity}[id=dia-10-corta-02,type=short,weight=1]

Completa: p\(\to\)q, p, por tanto __.

\end{activity}

\begin{activity}[id=dia-10-corta-03,type=short,weight=1]

Escribe el esquema de modus tollens.

\end{activity}

\begin{activity}[id=dia-10-corta-04,type=short,weight=1]

Completa: p\(\to\)q, \(\neg\)q, por tanto __.

\end{activity}

\begin{activity}[id=dia-10-corta-05,type=short,weight=1]

Identifica modus ponens en un ejemplo propio.

\end{activity}

\begin{activity}[id=dia-10-corta-06,type=short,weight=1]

¿Qué regla permite concluir \(\neg\)p a partir de p\(\to\)q y \(\neg\)q?

\end{activity}

\begin{activity}[id=dia-10-desarrollo-01,type=long,weight=2]

Resuelve dos argumentos mediante modus ponens y explica cada paso.

\end{activity}

\begin{activity}[id=dia-10-desarrollo-02,type=long,weight=2]

Construye un argumento válido por modus tollens e identifica claramente sus premisas y conclusión.

\end{activity}

\begin{activity}[id=dia-10-opcion-01,type=choice,weight=1]

De p\(\to\)q y p se infiere:

\option{a}{q}

\option{b}{\(\neg\)q}

\option{c}{\(\neg\)p}

\correct{a}

\end{activity}

\begin{activity}[id=dia-10-opcion-02,type=choice,weight=1]

De p\(\to\)q y \(\neg\)q se infiere:

\option{a}{p}

\option{b}{\(\neg\)p}

\option{c}{q}

\correct{b}

\end{activity}

\begin{deliverable}[id=entrega-dia-10,type=upload,weight=3]

Entrega un PDF con tres argumentos formalizados. Identifica premisas, conclusión y regla de inferencia en cada uno.

\end{deliverable}

\section{Día 11: Reglas de inferencia II}

\begin{lesson}[id=dia-11-lectura,weight=0]

Reconocer silogismo hipotético, disyuntivo y errores frecuentes.

\end{lesson}

\begin{activity}[id=dia-11-corta-01,type=short,weight=1]

Escribe el esquema del silogismo hipotético.

\end{activity}

\begin{activity}[id=dia-11-corta-02,type=short,weight=1]

Completa: p\(\lor\)q, \(\neg\)p, luego __.

\end{activity}

\begin{activity}[id=dia-11-corta-03,type=short,weight=1]

¿Qué se concluye de p\(\to\)q y q\(\to\)r?

\end{activity}

\begin{activity}[id=dia-11-corta-04,type=short,weight=1]

Identifica la regla: p\(\lor\)q, \(\neg\)q, por tanto p.

\end{activity}

\begin{activity}[id=dia-11-corta-05,type=short,weight=1]

¿Es válido afirmar p a partir de p\(\to\)q y q?

\end{activity}

\begin{activity}[id=dia-11-corta-06,type=short,weight=1]

Nombra el error de afirmar el consecuente.

\end{activity}

\begin{activity}[id=dia-11-desarrollo-01,type=long,weight=2]

Analiza un argumento que use silogismo disyuntivo. Explica por qué la conclusión se sigue.

\end{activity}

\begin{activity}[id=dia-11-desarrollo-02,type=long,weight=2]

Inventa un ejemplo de afirmación del consecuente y explica por qué no constituye una inferencia válida.

\end{activity}

\begin{activity}[id=dia-11-opcion-01,type=choice,weight=1]

De p\(\lor\)q y \(\neg\)p se obtiene:

\option{a}{\(\neg\)q}

\option{b}{q}

\option{c}{p}

\correct{b}

\end{activity}

\begin{activity}[id=dia-11-opcion-02,type=choice,weight=1]

De p\(\to\)q y q no se puede concluir necesariamente:

\option{a}{p}

\option{b}{q}

\option{c}{p\(\to\)q}

\correct{a}

\end{activity}

\section{Día 12: Axiomas y sistemas formales}

\begin{lesson}[id=dia-12-lectura,weight=0]

Comprender axiomas, definiciones, teoremas y reglas de inferencia.

\end{lesson}

\begin{activity}[id=dia-12-corta-01,type=short,weight=1]

Define axioma en una frase.

\end{activity}

\begin{activity}[id=dia-12-corta-02,type=short,weight=1]

¿Qué función cumple una definición?

\end{activity}

\begin{activity}[id=dia-12-corta-03,type=short,weight=1]

¿Qué es un teorema?

\end{activity}

\begin{activity}[id=dia-12-corta-04,type=short,weight=1]

¿Un axioma se demuestra dentro del sistema que lo adopta?

\end{activity}

\begin{activity}[id=dia-12-corta-05,type=short,weight=1]

Escribe un ejemplo de regla de inferencia.

\end{activity}

\begin{activity}[id=dia-12-corta-06,type=short,weight=1]

Ordena: axiomas, demostración, teorema (según su relación lógica).

\end{activity}

\begin{activity}[id=dia-12-desarrollo-01,type=long,weight=2]

Explica la diferencia entre aceptar un axioma y demostrar un teorema. Usa un ejemplo sencillo.

\end{activity}

\begin{activity}[id=dia-12-desarrollo-02,type=long,weight=2]

Describe cómo un sistema formal obtiene resultados a partir de axiomas, definiciones y reglas de inferencia.

\end{activity}

\begin{activity}[id=dia-12-opcion-01,type=choice,weight=1]

Un axioma es:

\option{a}{Una premisa adoptada como punto de partida}

\option{b}{Una conclusión ya demostrada}

\option{c}{Una operación aritmética}

\correct{a}

\end{activity}

\begin{activity}[id=dia-12-opcion-02,type=choice,weight=1]

Un teorema es una proposición:

\option{a}{Aceptada sin razones}

\option{b}{Demostrada a partir de supuestos}

\option{c}{Siempre falsa}

\correct{b}

\end{activity}

\section{Día 13: Demostraciones: estructura conceptual}

\begin{lesson}[id=dia-13-lectura,weight=0]

Identificar hipótesis, tesis y cadena de razones.

\end{lesson}

\begin{activity}[id=dia-13-corta-01,type=short,weight=1]

Señala la hipótesis en: si n es par, entonces n\(^{2}\) es par.

\end{activity}

\begin{activity}[id=dia-13-corta-02,type=short,weight=1]

Señala la tesis de la afirmación anterior.

\end{activity}

\begin{activity}[id=dia-13-corta-03,type=short,weight=1]

¿Qué significa demostrar una afirmación?

\end{activity}

\begin{activity}[id=dia-13-corta-04,type=short,weight=1]

¿Una lista de ejemplos prueba una afirmación universal?

\end{activity}

\begin{activity}[id=dia-13-corta-05,type=short,weight=1]

Escribe una afirmación matemática con hipótesis y conclusión.

\end{activity}

\begin{activity}[id=dia-13-corta-06,type=short,weight=1]

Identifica qué debe justificarse en una demostración.

\end{activity}

\begin{activity}[id=dia-13-desarrollo-01,type=long,weight=2]

Explica con tus palabras la estructura «supuestos \(\to\) razones \(\to\) conclusión». Incluye un ejemplo.

\end{activity}

\begin{activity}[id=dia-13-desarrollo-02,type=long,weight=2]

Analiza una demostración breve de tu elección: identifica hipótesis, pasos justificados y conclusión. Si falta una justificación, señálala.

\end{activity}

\begin{activity}[id=dia-13-opcion-01,type=choice,weight=1]

Una demostración matemática busca:

\option{a}{Convencer por autoridad}

\option{b}{Justificar lógicamente una afirmación}

\option{c}{Acumular ejemplos}

\correct{b}

\end{activity}

\begin{activity}[id=dia-13-opcion-02,type=choice,weight=1]

Probar varios casos particulares demuestra siempre un enunciado universal:

\option{a}{Sí}

\option{b}{No}

\option{c}{Solo si son tres}

\correct{b}

\end{activity}

\section{Día 14: Introducción a conjuntos}

\begin{lesson}[id=dia-14-lectura,weight=0]

Comprender conjunto, elemento, pertenencia y descripción.

\end{lesson}

\begin{activity}[id=dia-14-corta-01,type=short,weight=1]

Escribe por extensión el conjunto de vocales.

\end{activity}

\begin{activity}[id=dia-14-corta-02,type=short,weight=1]

Escribe por extensión los naturales menores que 5 (usa la convención indicada por tu curso: comienza en 0).

\end{activity}

\begin{activity}[id=dia-14-corta-03,type=short,weight=1]

Expresa 3\(\in\){1,2,3,4} con palabras.

\end{activity}

\begin{activity}[id=dia-14-corta-04,type=short,weight=1]

Decide si a\(\in\){a,b,c}.

\end{activity}

\begin{activity}[id=dia-14-corta-05,type=short,weight=1]

Escribe un conjunto de tres objetos matemáticos.

\end{activity}

\begin{activity}[id=dia-14-corta-06,type=short,weight=1]

Distingue conjunto y elemento con un ejemplo.

\end{activity}

\begin{activity}[id=dia-14-desarrollo-01,type=long,weight=2]

Explica la diferencia entre un objeto y el conjunto que lo contiene. Usa \(\in\) y \(\notin\) correctamente.

\end{activity}

\begin{activity}[id=dia-14-desarrollo-02,type=long,weight=2]

Describe el mismo conjunto por extensión y por comprensión, y verifica que ambas descripciones coincidan.

\end{activity}

\begin{activity}[id=dia-14-opcion-01,type=choice,weight=1]

El símbolo \(\in\) significa:

\option{a}{Está contenido como subconjunto}

\option{b}{Pertenece a}

\option{c}{Es igual a}

\correct{b}

\end{activity}

\begin{activity}[id=dia-14-opcion-02,type=choice,weight=1]

Si A={1,2,3}, entonces 4\(\in\)A es:

\option{a}{Verdadero}

\option{b}{Falso}

\option{c}{Una definición}

\correct{b}

\end{activity}

\section{Día 15: Notación y subconjuntos}

\begin{lesson}[id=dia-15-lectura,weight=0]

Usar inclusión, igualdad y conjunto vacío.

\end{lesson}

\begin{activity}[id=dia-15-corta-01,type=short,weight=1]

Escribe \(\varnothing\) con palabras.

\end{activity}

\begin{activity}[id=dia-15-corta-02,type=short,weight=1]

Decide si {1,2}\(\subseteq\){1,2,3}.

\end{activity}

\begin{activity}[id=dia-15-corta-03,type=short,weight=1]

Decide si {1,4}\(\subseteq\){1,2,3}.

\end{activity}

\begin{activity}[id=dia-15-corta-04,type=short,weight=1]

¿El conjunto vacío es subconjunto de todo conjunto?

\end{activity}

\begin{activity}[id=dia-15-corta-05,type=short,weight=1]

Escribe un subconjunto propio de {a,b,c}.

\end{activity}

\begin{activity}[id=dia-15-corta-06,type=short,weight=1]

Expresa que A y B tienen los mismos elementos.

\end{activity}

\begin{activity}[id=dia-15-desarrollo-01,type=long,weight=2]

Explica la diferencia entre x\(\in\)A y {x}\(\subseteq\)A mediante ejemplos.

\end{activity}

\begin{activity}[id=dia-15-desarrollo-02,type=long,weight=2]

Demuestra por definición que {1,2}\(\subseteq\){1,2,3}. Explica qué debes verificar.

\end{activity}

\begin{activity}[id=dia-15-opcion-01,type=choice,weight=1]

A\(\subseteq\)B significa:

\option{a}{Todo elemento de A pertenece a B}

\option{b}{A es elemento de B}

\option{c}{A y B son disjuntos}

\correct{a}

\end{activity}

\begin{activity}[id=dia-15-opcion-02,type=choice,weight=1]

El conjunto vacío:

\option{a}{No es subconjunto de ningún conjunto}

\option{b}{Es subconjunto de todo conjunto}

\option{c}{Tiene un elemento}

\correct{b}

\end{activity}

\begin{deliverable}[id=entrega-dia-15,type=upload,weight=3]

Entrega una hoja con definiciones de conjunto, elemento, pertenencia, subconjunto y conjunto universal, acompañadas de ejemplos correctos.

\end{deliverable}

\section{Día 16: Conjunto universal y complementario}

\begin{lesson}[id=dia-16-lectura,weight=0]

Definir universo de trabajo y complemento relativo.

\end{lesson}

\begin{activity}[id=dia-16-corta-01,type=short,weight=1]

¿Qué representa el conjunto universal U?

\end{activity}

\begin{activity}[id=dia-16-corta-02,type=short,weight=1]

Si U={1,2,3,4} y A={1,3}, escribe \(A^{c}\).

\end{activity}

\begin{activity}[id=dia-16-corta-03,type=short,weight=1]

¿El complemento depende de U?

\end{activity}

\begin{activity}[id=dia-16-corta-04,type=short,weight=1]

Si A=U, determina \(A^{c}\).

\end{activity}

\begin{activity}[id=dia-16-corta-05,type=short,weight=1]

Si A=\(\varnothing\), determina \(A^{c}\) respecto de U.

\end{activity}

\begin{activity}[id=dia-16-corta-06,type=short,weight=1]

Escribe un universo adecuado para el conjunto de vocales.

\end{activity}

\begin{activity}[id=dia-16-desarrollo-01,type=long,weight=2]

Explica por qué el complemento de un conjunto solo queda determinado cuando se especifica el universo.

\end{activity}

\begin{activity}[id=dia-16-desarrollo-02,type=long,weight=2]

Para U={1,2,3,4,5,6} y A={2,4,6}, calcula \(A^{c}\) y verifica que A\(\cup\)\(A^{c}\)=U y A\(\cap\)\(A^{c}\)=\(\varnothing\).

\end{activity}

\begin{activity}[id=dia-16-opcion-01,type=choice,weight=1]

El complemento \(A^{c}\) se define respecto de:

\option{a}{Un universo U}

\option{b}{Cualquier número}

\option{c}{El conjunto potencia}

\correct{a}

\end{activity}

\begin{activity}[id=dia-16-opcion-02,type=choice,weight=1]

Si A=U, entonces \(A^{c}\) es:

\option{a}{U}

\option{b}{\(\varnothing\)}

\option{c}{A}

\correct{b}

\end{activity}

\section{Día 17: Unión e intersección}

\begin{lesson}[id=dia-17-lectura,weight=0]

Calcular unión e intersección de conjuntos.

\end{lesson}

\begin{activity}[id=dia-17-corta-01,type=short,weight=1]

Calcula {1,2}\(\cup\){2,3}.

\end{activity}

\begin{activity}[id=dia-17-corta-02,type=short,weight=1]

Calcula {1,2}\(\cap\){2,3}.

\end{activity}

\begin{activity}[id=dia-17-corta-03,type=short,weight=1]

Calcula {a,b}\(\cup\){c}.

\end{activity}

\begin{activity}[id=dia-17-corta-04,type=short,weight=1]

Calcula {a,b}\(\cap\){c}.

\end{activity}

\begin{activity}[id=dia-17-corta-05,type=short,weight=1]

¿La unión conserva elementos repetidos?

\end{activity}

\begin{activity}[id=dia-17-corta-06,type=short,weight=1]

Interpreta A\(\cap\)B en palabras.

\end{activity}

\begin{activity}[id=dia-17-desarrollo-01,type=long,weight=2]

Resuelve unión e intersección de tres pares de conjuntos y explica el criterio de pertenencia en cada operación.

\end{activity}

\begin{activity}[id=dia-17-desarrollo-02,type=long,weight=2]

Da un ejemplo de conjuntos disjuntos y otro de conjuntos con intersección no vacía. Justifica ambos.

\end{activity}

\begin{activity}[id=dia-17-opcion-01,type=choice,weight=1]

La unión A\(\cup\)B contiene:

\option{a}{Los elementos de A o de B, o de ambos}

\option{b}{Solo los comunes}

\option{c}{Solo los de A}

\correct{a}

\end{activity}

\begin{activity}[id=dia-17-opcion-02,type=choice,weight=1]

La intersección A\(\cap\)B contiene:

\option{a}{Todos los elementos distintos}

\option{b}{Los elementos comunes}

\option{c}{Los elementos externos}

\correct{b}

\end{activity}

\section{Día 18: Diferencia y diferencia simétrica}

\begin{lesson}[id=dia-18-lectura,weight=0]

Operar con diferencia de conjuntos.

\end{lesson}

\begin{activity}[id=dia-18-corta-01,type=short,weight=1]

Calcula {1,2,3}\{2}.

\end{activity}

\begin{activity}[id=dia-18-corta-02,type=short,weight=1]

Calcula {a,b}\{b,c}.

\end{activity}

\begin{activity}[id=dia-18-corta-03,type=short,weight=1]

Calcula {1,2}\(\triangle\){2,3}.

\end{activity}

\begin{activity}[id=dia-18-corta-04,type=short,weight=1]

¿La diferencia A\B siempre coincide con B\A?

\end{activity}

\begin{activity}[id=dia-18-corta-05,type=short,weight=1]

Describe A\B en palabras.

\end{activity}

\begin{activity}[id=dia-18-corta-06,type=short,weight=1]

¿Qué elementos excluye la diferencia simétrica?

\end{activity}

\begin{activity}[id=dia-18-desarrollo-01,type=long,weight=2]

Calcula A\B y B\A para A={1,2,3} y B={2,3,4}. Explica la diferencia.

\end{activity}

\begin{activity}[id=dia-18-desarrollo-02,type=long,weight=2]

Construye un ejemplo de diferencia simétrica y verifica que contiene los elementos que pertenecen exactamente a uno de los dos conjuntos.

\end{activity}

\begin{activity}[id=dia-18-opcion-01,type=choice,weight=1]

A\B contiene elementos que:

\option{a}{Están en A y no en B}

\option{b}{Están en ambos}

\option{c}{No están en ninguno}

\correct{a}

\end{activity}

\begin{activity}[id=dia-18-opcion-02,type=choice,weight=1]

La diferencia simétrica contiene elementos que pertenecen:

\option{a}{A ambos conjuntos}

\option{b}{Exactamente a uno de los dos}

\option{c}{A ninguno}

\correct{b}

\end{activity}

\section{Día 19: Leyes de conjuntos}

\begin{lesson}[id=dia-19-lectura,weight=0]

Aplicar conmutatividad, asociatividad, distributividad e identidad.

\end{lesson}

\begin{activity}[id=dia-19-corta-01,type=short,weight=1]

Escribe la ley conmutativa de la unión.

\end{activity}

\begin{activity}[id=dia-19-corta-02,type=short,weight=1]

Escribe la ley conmutativa de la intersección.

\end{activity}

\begin{activity}[id=dia-19-corta-03,type=short,weight=1]

Completa: A\(\cup\)\(\varnothing\)=__.

\end{activity}

\begin{activity}[id=dia-19-corta-04,type=short,weight=1]

Completa: A\(\cap\)U=__.

\end{activity}

\begin{activity}[id=dia-19-corta-05,type=short,weight=1]

Escribe la ley distributiva de \(\cap\) respecto de \(\cup\).

\end{activity}

\begin{activity}[id=dia-19-corta-06,type=short,weight=1]

¿La unión es asociativa?

\end{activity}

\begin{activity}[id=dia-19-desarrollo-01,type=long,weight=2]

Verifica una ley de conjuntos con un ejemplo concreto y explica por qué un ejemplo no sustituye una demostración general.

\end{activity}

\begin{activity}[id=dia-19-desarrollo-02,type=long,weight=2]

Enuncia las leyes de identidad, dominación y complemento con notación correcta.

\end{activity}

\begin{activity}[id=dia-19-opcion-01,type=choice,weight=1]

A\(\cup\)\(\varnothing\) es igual a:

\option{a}{\(\varnothing\)}

\option{b}{A}

\option{c}{U}

\correct{b}

\end{activity}

\begin{activity}[id=dia-19-opcion-02,type=choice,weight=1]

A\(\cap\)U es igual a:

\option{a}{A}

\option{b}{\(\varnothing\)}

\option{c}{U}

\correct{a}

\end{activity}

\section{Día 20: Diagramas de Venn}

\begin{lesson}[id=dia-20-lectura,weight=0]

Representar operaciones y regiones de conjuntos.

\end{lesson}

\begin{activity}[id=dia-20-corta-01,type=short,weight=1]

Sombrea conceptualmente A\(\cup\)B en un diagrama de dos conjuntos.

\end{activity}

\begin{activity}[id=dia-20-corta-02,type=short,weight=1]

Describe la región A\(\cap\)B.

\end{activity}

\begin{activity}[id=dia-20-corta-03,type=short,weight=1]

Describe la región A\B.

\end{activity}

\begin{activity}[id=dia-20-corta-04,type=short,weight=1]

¿Qué región representa (A\(\cup\)B)ᶜ?

\end{activity}

\begin{activity}[id=dia-20-corta-05,type=short,weight=1]

¿Qué representa A\(\cap\)B=\(\varnothing\) en un diagrama?

\end{activity}

\begin{activity}[id=dia-20-corta-06,type=short,weight=1]

Dibuja dos conjuntos dentro de U.

\end{activity}

\begin{activity}[id=dia-20-desarrollo-01,type=long,weight=2]

Entrega un diagrama de Venn rotulado para U, A y B, y representa por separado unión, intersección y complemento.

\end{activity}

\begin{activity}[id=dia-20-desarrollo-02,type=long,weight=2]

Explica cómo un diagrama de Venn ayuda a interpretar operaciones, e indica una limitación de usar un dibujo como demostración general.

\end{activity}

\begin{activity}[id=dia-20-opcion-01,type=choice,weight=1]

La región común a A y B representa:

\option{a}{A\(\cup\)B}

\option{b}{A\(\cap\)B}

\option{c}{A\B}

\correct{b}

\end{activity}

\begin{activity}[id=dia-20-opcion-02,type=choice,weight=1]

El rectángulo exterior suele representar:

\option{a}{U}

\option{b}{\(\varnothing\)}

\option{c}{A\(\cap\)B}

\correct{a}

\end{activity}

\begin{deliverable}[id=entrega-dia-20,type=upload,weight=3]

Entrega un PDF con ejercicios resueltos de unión, intersección, diferencia y complemento. Indica explícitamente el universo.

\end{deliverable}

\section{Día 21: Axiomas básicos de conjuntos: extensionalidad}

\begin{lesson}[id=dia-21-lectura,weight=0]

Comprender igualdad de conjuntos por sus elementos.

\end{lesson}

\begin{activity}[id=dia-21-corta-01,type=short,weight=1]

Enuncia el principio de extensionalidad.

\end{activity}

\begin{activity}[id=dia-21-corta-02,type=short,weight=1]

¿Son iguales {1,2} y {2,1}?

\end{activity}

\begin{activity}[id=dia-21-corta-03,type=short,weight=1]

¿Son iguales {1,1,2} y {1,2}?

\end{activity}

\begin{activity}[id=dia-21-corta-04,type=short,weight=1]

Si A\(\subseteq\)B y B\(\subseteq\)A, ¿qué se concluye?

\end{activity}

\begin{activity}[id=dia-21-corta-05,type=short,weight=1]

¿El orden de escritura cambia un conjunto?

\end{activity}

\begin{activity}[id=dia-21-corta-06,type=short,weight=1]

¿La repetición de elementos cambia un conjunto?

\end{activity}

\begin{activity}[id=dia-21-desarrollo-01,type=long,weight=2]

Explica por qué dos conjuntos son iguales exactamente cuando tienen los mismos elementos.

\end{activity}

\begin{activity}[id=dia-21-desarrollo-02,type=long,weight=2]

Demuestra que {a,b}={b,a} usando extensionalidad o doble inclusión.

\end{activity}

\begin{activity}[id=dia-21-opcion-01,type=choice,weight=1]

Dos conjuntos son iguales cuando:

\option{a}{Tienen el mismo nombre}

\option{b}{Tienen exactamente los mismos elementos}

\option{c}{Se escriben en el mismo orden}

\correct{b}

\end{activity}

\begin{activity}[id=dia-21-opcion-02,type=choice,weight=1]

Si A\(\subseteq\)B y B\(\subseteq\)A, entonces:

\option{a}{A=B}

\option{b}{A\(\cap\)B=\(\varnothing\)}

\option{c}{A\(\in\)B}

\correct{a}

\end{activity}

\section{Día 22: Axiomas básicos: pares, unión y vacío}

\begin{lesson}[id=dia-22-lectura,weight=0]

Conocer construcciones elementales de conjuntos.

\end{lesson}

\begin{activity}[id=dia-22-corta-01,type=short,weight=1]

¿Qué conjunto forma el axioma del par a partir de a y b?

\end{activity}

\begin{activity}[id=dia-22-corta-02,type=short,weight=1]

¿Qué reúne el axioma de unión?

\end{activity}

\begin{activity}[id=dia-22-corta-03,type=short,weight=1]

Escribe un conjunto que contenga a y b como elementos.

\end{activity}

\begin{activity}[id=dia-22-corta-04,type=short,weight=1]

¿Qué caracteriza al conjunto vacío?

\end{activity}

\begin{activity}[id=dia-22-corta-05,type=short,weight=1]

¿Puede haber más de una forma de describir el mismo conjunto?

\end{activity}

\begin{activity}[id=dia-22-corta-06,type=short,weight=1]

¿Qué significa que una colección esté bien definida?

\end{activity}

\begin{activity}[id=dia-22-desarrollo-01,type=long,weight=2]

Explica intuitivamente cómo el axioma del par y el axioma de unión permiten construir conjuntos.

\end{activity}

\begin{activity}[id=dia-22-desarrollo-02,type=long,weight=2]

Da ejemplos de un conjunto de dos elementos y de la unión de dos conjuntos. Describe qué elementos aparecen en cada construcción.

\end{activity}

\begin{activity}[id=dia-22-opcion-01,type=choice,weight=1]

El axioma del par permite formar:

\option{a}{Un conjunto con a y b como elementos}

\option{b}{Un número primo}

\option{c}{Una función}

\correct{a}

\end{activity}

\begin{activity}[id=dia-22-opcion-02,type=choice,weight=1]

El conjunto vacío se caracteriza por:

\option{a}{Tener un elemento}

\option{b}{No tener elementos}

\option{c}{Tener todos los elementos}

\correct{b}

\end{activity}

\section{Día 23: Axiomas básicos: separación y potencia}

\begin{lesson}[id=dia-23-lectura,weight=0]

Entender la formación de subconjuntos y el conjunto potencia.

\end{lesson}

\begin{activity}[id=dia-23-corta-01,type=short,weight=1]

¿Qué reúne P(A)?

\end{activity}

\begin{activity}[id=dia-23-corta-02,type=short,weight=1]

Escribe P({a}).

\end{activity}

\begin{activity}[id=dia-23-corta-03,type=short,weight=1]

¿Cuántos subconjuntos tiene un conjunto de dos elementos?

\end{activity}

\begin{activity}[id=dia-23-corta-04,type=short,weight=1]

¿\(\varnothing\) pertenece a P(A) para todo A?

\end{activity}

\begin{activity}[id=dia-23-corta-05,type=short,weight=1]

¿A pertenece a P(A)?

\end{activity}

\begin{activity}[id=dia-23-corta-06,type=short,weight=1]

Si A tiene tres elementos, ¿cuántos elementos tiene P(A)?

\end{activity}

\begin{activity}[id=dia-23-desarrollo-01,type=long,weight=2]

Explica la idea del axioma de separación: formar un subconjunto mediante una propiedad dentro de un conjunto dado.

\end{activity}

\begin{activity}[id=dia-23-desarrollo-02,type=long,weight=2]

Construye el conjunto potencia de {a,b} y verifica que cada uno de sus elementos sea un subconjunto del conjunto original.

\end{activity}

\begin{activity}[id=dia-23-opcion-01,type=choice,weight=1]

P(A) es:

\option{a}{El conjunto de todos los subconjuntos de A}

\option{b}{El complemento de A}

\option{c}{La unión de A consigo mismo}

\correct{a}

\end{activity}

\begin{activity}[id=dia-23-opcion-02,type=choice,weight=1]

Si |A|=3, entonces |P(A)| es:

\option{a}{3}

\option{b}{6}

\option{c}{8}

\correct{c}

\end{activity}

\section{Día 24: Conjunto potencia}

\begin{lesson}[id=dia-24-lectura,weight=0]

Enumerar conjuntos potencia y reconocer su cardinalidad finita.

\end{lesson}

\begin{activity}[id=dia-24-corta-01,type=short,weight=1]

Escribe P(\(\varnothing\)).

\end{activity}

\begin{activity}[id=dia-24-corta-02,type=short,weight=1]

Escribe P({1}).

\end{activity}

\begin{activity}[id=dia-24-corta-03,type=short,weight=1]

Enumera P({a,b}).

\end{activity}

\begin{activity}[id=dia-24-corta-04,type=short,weight=1]

¿Cuántos elementos tiene P({1,2,3,4})?

\end{activity}

\begin{activity}[id=dia-24-corta-05,type=short,weight=1]

¿Todo elemento de P(A) es un conjunto?

\end{activity}

\begin{activity}[id=dia-24-corta-06,type=short,weight=1]

¿A es siempre elemento de P(A)?

\end{activity}

\begin{activity}[id=dia-24-desarrollo-01,type=long,weight=2]

Enumera P({1,2,3}) sin omitir subconjuntos. Cuenta y verifica el resultado.

\end{activity}

\begin{activity}[id=dia-24-desarrollo-02,type=long,weight=2]

Explica por qué un conjunto finito con n elementos tiene 2\(^{n}\) subconjuntos, usando el caso de tres elementos.

\end{activity}

\begin{activity}[id=dia-24-opcion-01,type=choice,weight=1]

P({a,b}) contiene:

\option{a}{2 subconjuntos}

\option{b}{3 subconjuntos}

\option{c}{4 subconjuntos}

\correct{c}

\end{activity}

\begin{activity}[id=dia-24-opcion-02,type=choice,weight=1]

Si A tiene n elementos, |P(A)| es:

\option{a}{n\(^{2}\)}

\option{b}{2\(^{n}\)}

\option{c}{2n}

\correct{b}

\end{activity}

\section{Día 25: Cuantificadores y conjuntos}

\begin{lesson}[id=dia-25-lectura,weight=0]

Interpretar «para todo» y «existe» en propiedades de conjuntos.

\end{lesson}

\begin{activity}[id=dia-25-corta-01,type=short,weight=1]

Escribe el símbolo de «para todo».

\end{activity}

\begin{activity}[id=dia-25-corta-02,type=short,weight=1]

Escribe el símbolo de «existe».

\end{activity}

\begin{activity}[id=dia-25-corta-03,type=short,weight=1]

Traduce: todo x\(\in\)A también pertenece a B.

\end{activity}

\begin{activity}[id=dia-25-corta-04,type=short,weight=1]

Traduce: existe x\(\in\)A que no pertenece a B.

\end{activity}

\begin{activity}[id=dia-25-corta-05,type=short,weight=1]

¿Qué expresa x\(\in\)A \(\land\) x\(\in\)B?

\end{activity}

\begin{activity}[id=dia-25-corta-06,type=short,weight=1]

Escribe la condición de A\(\subseteq\)B con cuantificadores.

\end{activity}

\begin{activity}[id=dia-25-desarrollo-01,type=long,weight=2]

Formula con cuantificadores la igualdad A=B mediante doble inclusión.

\end{activity}

\begin{activity}[id=dia-25-desarrollo-02,type=long,weight=2]

Explica la diferencia entre \(\forall\)x\(\in\)A y \(\exists\)x\(\in\)A mediante dos afirmaciones sobre un conjunto concreto.

\end{activity}

\begin{activity}[id=dia-25-opcion-01,type=choice,weight=1]

\(\forall\) se lee:

\option{a}{Existe}

\option{b}{Para todo}

\option{c}{No pertenece}

\correct{b}

\end{activity}

\begin{activity}[id=dia-25-opcion-02,type=choice,weight=1]

A\(\subseteq\)B equivale a:

\option{a}{\(\forall\)x(x\(\in\)A\(\to\)x\(\in\)B)}

\option{b}{\(\exists\)x(x\(\in\)A\(\land\)x\(\notin\)B)}

\option{c}{A\(\in\)B}

\correct{a}

\end{activity}

\begin{deliverable}[id=entrega-dia-25,type=upload,weight=3]

Entrega un PDF que enumere el conjunto potencia de dos conjuntos finitos y explique el número de subconjuntos.

\end{deliverable}

\section{Día 26: Demostraciones por inclusión}

\begin{lesson}[id=dia-26-lectura,weight=0]

Probar inclusiones siguiendo la definición.

\end{lesson}

\begin{activity}[id=dia-26-corta-01,type=short,weight=1]

¿Qué debes suponer para probar A\(\subseteq\)B?

\end{activity}

\begin{activity}[id=dia-26-corta-02,type=short,weight=1]

¿Qué debes demostrar sobre un elemento x?

\end{activity}

\begin{activity}[id=dia-26-corta-03,type=short,weight=1]

Completa: sea x\(\in\)A; entonces... (objetivo: x\(\in\)B).

\end{activity}

\begin{activity}[id=dia-26-corta-04,type=short,weight=1]

¿Probar A\(\subseteq\)B basta para concluir A=B?

\end{activity}

\begin{activity}[id=dia-26-corta-05,type=short,weight=1]

¿Qué segunda inclusión se requiere para probar igualdad?

\end{activity}

\begin{activity}[id=dia-26-corta-06,type=short,weight=1]

Escribe el esquema de una prueba por inclusión.

\end{activity}

\begin{activity}[id=dia-26-desarrollo-01,type=long,weight=2]

Demuestra que A\(\cap\)B\(\subseteq\)A usando la definición de intersección.

\end{activity}

\begin{activity}[id=dia-26-desarrollo-02,type=long,weight=2]

Demuestra que A\(\subseteq\)A\(\cup\)B usando la definición de unión y organiza la prueba en pasos.

\end{activity}

\begin{activity}[id=dia-26-opcion-01,type=choice,weight=1]

Para probar A\(\subseteq\)B se toma:

\option{a}{Un elemento arbitrario de A}

\option{b}{Un elemento externo}

\option{c}{Solo el conjunto vacío}

\correct{a}

\end{activity}

\begin{activity}[id=dia-26-opcion-02,type=choice,weight=1]

Para demostrar A=B basta probar:

\option{a}{A\(\subseteq\)B}

\option{b}{B\(\subseteq\)A}

\option{c}{A\(\subseteq\)B y B\(\subseteq\)A}

\correct{c}

\end{activity}

\section{Día 27: Demostraciones de identidades de conjuntos}

\begin{lesson}[id=dia-27-lectura,weight=0]

Demostrar identidades con doble inclusión.

\end{lesson}

\begin{activity}[id=dia-27-corta-01,type=short,weight=1]

¿Qué estrategia sirve para demostrar A=B?

\end{activity}

\begin{activity}[id=dia-27-corta-02,type=short,weight=1]

Para probar A\(\cup\)B=B\(\cup\)A, ¿qué debes tomar arbitrariamente?

\end{activity}

\begin{activity}[id=dia-27-corta-03,type=short,weight=1]

¿Qué significa x\(\in\)A\(\cup\)B?

\end{activity}

\begin{activity}[id=dia-27-corta-04,type=short,weight=1]

¿Qué significa x\(\in\)A\(\cap\)B?

\end{activity}

\begin{activity}[id=dia-27-corta-05,type=short,weight=1]

¿Qué dos direcciones se prueban en doble inclusión?

\end{activity}

\begin{activity}[id=dia-27-corta-06,type=short,weight=1]

¿Un diagrama por sí solo prueba una identidad general?

\end{activity}

\begin{activity}[id=dia-27-desarrollo-01,type=long,weight=2]

Demuestra A\(\cup\)B=B\(\cup\)A por doble inclusión.

\end{activity}

\begin{activity}[id=dia-27-desarrollo-02,type=long,weight=2]

Demuestra A\(\cap\)B=B\(\cap\)A por doble inclusión, justificando cada equivalencia.

\end{activity}

\begin{activity}[id=dia-27-opcion-01,type=choice,weight=1]

La doble inclusión consiste en probar:

\option{a}{A\(\subseteq\)B y B\(\subseteq\)A}

\option{b}{A\(\in\)B}

\option{c}{A\(\cap\)B=\(\varnothing\)}

\correct{a}

\end{activity}

\begin{activity}[id=dia-27-opcion-02,type=choice,weight=1]

x\(\in\)A\(\cap\)B significa:

\option{a}{x\(\in\)A o x\(\in\)B}

\option{b}{x\(\in\)A y x\(\in\)B}

\option{c}{x\(\notin\)A}

\correct{b}

\end{activity}

\section{Día 28: Repaso integrado de lógica}

\begin{lesson}[id=dia-28-lectura,weight=0]

Integrar proposiciones, tablas, equivalencias e inferencias.

\end{lesson}

\begin{activity}[id=dia-28-corta-01,type=short,weight=1]

Construye la negación de p\(\land\)q.

\end{activity}

\begin{activity}[id=dia-28-corta-02,type=short,weight=1]

¿Qué regla aplica a p\(\to\)q y p?

\end{activity}

\begin{activity}[id=dia-28-corta-03,type=short,weight=1]

Clasifica p\(\lor\)\(\neg\)p.

\end{activity}

\begin{activity}[id=dia-28-corta-04,type=short,weight=1]

¿Cuántas filas tiene una tabla de cuatro variables?

\end{activity}

\begin{activity}[id=dia-28-corta-05,type=short,weight=1]

Traduce «si p entonces q» a símbolos.

\end{activity}

\begin{activity}[id=dia-28-corta-06,type=short,weight=1]

Identifica la conclusión en un argumento de dos premisas.

\end{activity}

\begin{activity}[id=dia-28-desarrollo-01,type=long,weight=2]

Construye una tabla de verdad para (p\(\to\)q)\(\lor\)\(\neg\)q y clasifícala.

\end{activity}

\begin{activity}[id=dia-28-desarrollo-02,type=long,weight=2]

Analiza un argumento de tres líneas: formalízalo, identifica la regla de inferencia y determina si es válido.

\end{activity}

\begin{activity}[id=dia-28-opcion-01,type=choice,weight=1]

De p\(\to\)q y p se concluye:

\option{a}{q}

\option{b}{\(\neg\)q}

\option{c}{\(\neg\)p}

\correct{a}

\end{activity}

\begin{activity}[id=dia-28-opcion-02,type=choice,weight=1]

p\(\lor\)\(\neg\)p es:

\option{a}{Contradicción}

\option{b}{Contingencia}

\option{c}{Tautología}

\correct{c}

\end{activity}

\section{Día 29: Repaso integrado de conjuntos}

\begin{lesson}[id=dia-29-lectura,weight=0]

Integrar notación, operaciones, universo y potencia.

\end{lesson}

\begin{activity}[id=dia-29-corta-01,type=short,weight=1]

Calcula A\(\cup\)B para A={1,2}, B={2,3}.

\end{activity}

\begin{activity}[id=dia-29-corta-02,type=short,weight=1]

Calcula A\(\cap\)B para los mismos conjuntos.

\end{activity}

\begin{activity}[id=dia-29-corta-03,type=short,weight=1]

Si U={1,2,3,4}, calcula el complemento de {1,4}.

\end{activity}

\begin{activity}[id=dia-29-corta-04,type=short,weight=1]

Escribe P({x}).

\end{activity}

\begin{activity}[id=dia-29-corta-05,type=short,weight=1]

¿Qué significa A\(\subseteq\)B?

\end{activity}

\begin{activity}[id=dia-29-corta-06,type=short,weight=1]

¿Qué elemento pertenece siempre a P(A)?

\end{activity}

\begin{activity}[id=dia-29-desarrollo-01,type=long,weight=2]

Resuelve un ejercicio que combine unión, intersección y complemento respecto de un universo explícito.

\end{activity}

\begin{activity}[id=dia-29-desarrollo-02,type=long,weight=2]

Construye P({a,b}) y explica cómo se relaciona cada elemento del conjunto potencia con el conjunto original.

\end{activity}

\begin{activity}[id=dia-29-opcion-01,type=choice,weight=1]

Si A={1,2} y B={2,3}, A\(\cap\)B es:

\option{a}{{1,2,3}}

\option{b}{{2}}

\option{c}{\(\varnothing\)}

\correct{b}

\end{activity}

\begin{activity}[id=dia-29-opcion-02,type=choice,weight=1]

P({x}) es:

\option{a}{{\(\varnothing\),{x}}}

\option{b}{{{x}}}

\option{c}{{x}}

\correct{a}

\end{activity}

\section{Día 30: Evaluación de demostraciones}

\begin{lesson}[id=dia-30-lectura,weight=0]

Redactar pruebas breves y reconocer justificaciones insuficientes.

\end{lesson}

\begin{activity}[id=dia-30-corta-01,type=short,weight=1]

¿Qué se supone al iniciar una prueba de inclusión?

\end{activity}

\begin{activity}[id=dia-30-corta-02,type=short,weight=1]

¿Qué se debe probar al final?

\end{activity}

\begin{activity}[id=dia-30-corta-03,type=short,weight=1]

¿Qué función cumple una definición en una demostración?

\end{activity}

\begin{activity}[id=dia-30-corta-04,type=short,weight=1]

¿Por qué debe justificarse cada paso?

\end{activity}

\begin{activity}[id=dia-30-corta-05,type=short,weight=1]

¿Qué diferencia hay entre ejemplo y prueba?

\end{activity}

\begin{activity}[id=dia-30-corta-06,type=short,weight=1]

Escribe una conclusión matemática con «por tanto».

\end{activity}

\begin{activity}[id=dia-30-desarrollo-01,type=long,weight=2]

Redacta una demostración completa de una inclusión sencilla, citando la definición utilizada.

\end{activity}

\begin{activity}[id=dia-30-desarrollo-02,type=long,weight=2]

Revisa una demostración propia y marca hipótesis, argumento, justificación y conclusión. Corrige cualquier salto lógico.

\end{activity}

\begin{activity}[id=dia-30-opcion-01,type=choice,weight=1]

Una prueba debe justificar:

\option{a}{Solo el primer paso}

\option{b}{Cada inferencia relevante}

\option{c}{Solo el resultado}

\correct{b}

\end{activity}

\begin{activity}[id=dia-30-opcion-02,type=choice,weight=1]

Un contraejemplo sirve para:

\option{a}{Probar todo enunciado verdadero}

\option{b}{Refutar un enunciado universal}

\option{c}{Definir un conjunto}

\correct{b}

\end{activity}

\begin{deliverable}[id=entrega-dia-30,type=upload,weight=3]

Entrega un PDF con dos demostraciones por doble inclusión y un problema integrado de lógica y conjuntos. Escribe todos los pasos y justificaciones.

\end{deliverable}

\section{Día 31: Síntesis y evaluación final}

\begin{lesson}[id=dia-31-lectura,weight=0]

Consolidar los conceptos centrales del curso.

\end{lesson}

\begin{activity}[id=dia-31-corta-01,type=short,weight=1]

Define proposición.

\end{activity}

\begin{activity}[id=dia-31-corta-02,type=short,weight=1]

Define axioma.

\end{activity}

\begin{activity}[id=dia-31-corta-03,type=short,weight=1]

Escribe una regla de inferencia válida.

\end{activity}

\begin{activity}[id=dia-31-corta-04,type=short,weight=1]

¿Qué es una tautología?

\end{activity}

\begin{activity}[id=dia-31-corta-05,type=short,weight=1]

Define subconjunto.

\end{activity}

\begin{activity}[id=dia-31-corta-06,type=short,weight=1]

Define conjunto potencia.

\end{activity}

\begin{activity}[id=dia-31-desarrollo-01,type=long,weight=2]

Escribe una síntesis de una página como máximo que conecte lógica proposicional, inferencia, axiomas, conjuntos y demostración.

\end{activity}

\begin{activity}[id=dia-31-desarrollo-02,type=long,weight=2]

Resuelve un problema integrado: formaliza dos premisas, deriva una conclusión válida y después expresa una colección de objetos como conjunto, calcula una operación y justifica una inclusión.

\end{activity}

\begin{activity}[id=dia-31-opcion-01,type=choice,weight=1]

Una demostración es:

\option{a}{Una colección de ejemplos}

\option{b}{Una justificación lógica}

\option{c}{Una opinión}

\correct{b}

\end{activity}

\begin{activity}[id=dia-31-opcion-02,type=choice,weight=1]

El conjunto potencia P(A) contiene:

\option{a}{Los elementos de A únicamente}

\option{b}{Todos los subconjuntos de A}

\option{c}{Solo el conjunto vacío}

\correct{b}

\end{activity}

\section{Referencia breve}

\begin{lesson}[id=referencia-conectivos,weight=0]

\begin{table}[id=referencia-conectivos-tabla,header=teal]

\row{Símbolo}{Lectura}

\row{\(\neg p\)}{No \(p\)}

\row{\(p\land q\)}{\(p\) y \(q\)}

\row{\(p\lor q\)}{\(p\) o \(q\), inclusiva}

\row{\(p\to q\)}{Si \(p\), entonces \(q\)}

\row{\(p\leftrightarrow q\)}{\(p\) si y solo si \(q\)}

\row{\(p\equiv q\)}{\(p\) y \(q\) son lógicamente equivalentes}

\end{table}

\end{lesson}

\begin{lesson}[id=referencia-inferencias,weight=0]

\begin{table}[id=referencia-inferencias-tabla,header=teal]

\row{Regla}{Esquema}

\row{Modus ponens}{\(p\to q,\ p\ \therefore q\)}

\row{Modus tollens}{\(p\to q,\ \neg q\ \therefore\neg p\)}

\row{Silogismo hipotético}{\(p\to q,\ q\to r\ \therefore p\to r\)}

\row{Silogismo disyuntivo}{\(p\lor q,\ \neg p\ \therefore q\)}

\end{table}

\end{lesson}

\begin{lesson}[id=referencia-conjuntos,weight=0]

\begin{table}[id=referencia-conjuntos-tabla,header=teal]

\row{Símbolo}{Lectura}

\row{\(x\in A\)}{\(x\) pertenece a \(A\)}

\row{\(A\subseteq B\)}{Todo elemento de \(A\) pertenece a \(B\)}

\row{\(A\cup B\)}{Unión}

\row{\(A\cap B\)}{Intersección}

\row{\(A\setminus B\)}{Diferencia}

\row{\(A^{c}\)}{Complemento respecto de \(U\)}

\row{\(P(A)\)}{Conjunto de todos los subconjuntos de \(A\)}

\end{table}

\end{lesson}

\begin{lesson}[id=referencia-axiomas-conjuntos,weight=0]

En una presentación axiomática elemental pueden estudiarse, entre otros, extensionalidad, conjunto vacío, par, unión y separación; el conjunto potencia se introduce mediante el axioma correspondiente. La formulación exacta depende del sistema axiomático adoptado. En este curso se trabaja con su significado intuitivo y su uso en ejercicios, sin desarrollar un tratamiento formal completo de ZF.

\end{lesson}

\begin{lesson}[id=referencia-esquema-demostracion,weight=0]

Para demostrar \(A\subseteq B\): toma un elemento arbitrario \(x\in A\), usa las definiciones y demuestra que \(x\in B\). Para demostrar \(A=B\): prueba \(A\subseteq B\) y \(B\subseteq A\). Cada paso debe estar justificado.

\end{lesson}
